\section{The Appointed Texts}
\label{sec:appointed-texts}

The Latin below is that of the 1962 typical edition, \work{Missale Romanum}
(Typis Polyglottis Vaticanis), pp.~412--413, nos.~1689--1698, read on the page
images of the CMAA facsimile. Every element's wording already stands in the
Pustet \work{Missale Romanum} of Ratisbon, 1862 (pp.~352--354); the two books
differ only in spelling, pointing, the form of the citations, the length of one
printed conclusion, and the additional orations that the 1862 book prints under
the rubrics of its day and the 1962 book no longer carries. The English of the
scriptural elements is the Douay--Rheims Bible in Bishop Challoner's revision,
which translates the Clementine Vulgate at the Missal's own Vulgate numbering.
It is given at the canonical verses, so that wherever a chant departs from its
verse the English and the Latin part company; each such place is named beneath
the text. The English of the Introit's antiphon and of the three orations is
the anonymous translation printed in \work{The Roman Missal, Translated into the
English Language for the Use of the Laity} (Philadelphia: Eugene Cummiskey,
1861). Neither English is an approved liturgical translation of the 1962
Missal. Psalms are cited in the Missal's numbering; the Hebrew numbering
follows in parentheses at each psalm's first citation and in the loci below.

\subsection{Introit}
\label{proper-introit}

\properlocus{Antiphon, printed citation \latin{Dan.\ 3, 31, 29 et 35}; verse Ps 118:1 (Heb.\ 119). Missal p.~412, no.~1689.}

\begin{latinproper}
Omnia, quæ fecísti nobis, Dómine, in vero iudício fecísti, quia peccávimus
tibi, et mandátis tuis non obœdívimus: sed da glóriam nómini tuo, et fac
nobíscum secúndum multitúdinem misericórdiæ tuæ. \textnormal{Ps.\ 118, 1}
Beáti immaculáti in via: qui ámbulant in lege Dómini. ℣.~Glória Patri.
\end{latinproper}

\begin{englishwitness}{Cummiskey, The Roman Missal (1861), XX Sunday after Pentecost, Introit}
Whatever thou hast done to us, O Lord, thou hast done by a just judgment: for
we have sinned and disobeyed thy commandments: but glorify thy name, and deal
with us according to thy great mercy.
\end{englishwitness}

\begin{englishwitness}{Douay--Rheims (Challoner), Ps 118:1, the verse}
Blessed are the undefiled in the way, who walk in the law of the Lord.
\end{englishwitness}

\noindent The antiphon is not a continuous text of the Vulgate, and its words do
not stand at all the verses the Missal cites. Its first clause shares its words
with Dan 3:31, ``all that thou hast brought upon us, and every thing that thou
hast done to us, thou hast done in true judgment''; \latin{peccavimus} stands
at verse 29, ``For we have sinned, and committed iniquity''; \latin{mandatis
tuis non oboedivimus} has no word-for-word counterpart, verse 30 reading ``we
have not hearkened to thy commandments, nor have we observed''; and the closing
petition gathers words of verses 43 and 42, ``give glory to thy name, O Lord''
and ``deal with us according to thy meekness, and according to the multitude of
thy mercies''. Verse 35, which the Missal cites, asks God not to take away his
mercy ``for the sake of Abraham, thy beloved, and Isaac, thy servant, and
Israel, thy holy one''. The Douay--Rheims therefore cannot carry the antiphon,
and the 1861 English is that book's own rendering of the Missal's Latin. It
leaves out \latin{tibi}, ``against thee'', after ``we have sinned'', and it
renders \latin{da gloriam nomini tuo} as ``glorify thy name'' and
\latin{secundum multitudinem misericordiae tuae} as ``according to thy great
mercy''. The Vulgate counts the psalm's title, \latin{Alleluia}, and the
letter \latin{Aleph} inside its first verse; the Missal's verse begins after
them. After the \latin{Gloria Patri} the antiphon is repeated (RGMR 427).

\subsection{Collect}
\label{proper-collect}

\properlocus{\latin{Oratio}. Missal p.~412, no.~1690.}

\begin{latinproper}
Largíre, quǽsumus, Dómine, fidélibus tuis indulgéntiam placátus et pacem: ut
páriter ab ómnibus mundéntur offénsis, et secúra tibi mente desérviant. Per
Dóminum.
\end{latinproper}

\begin{englishwitness}{Cummiskey, The Roman Missal (1861), XX Sunday after Pentecost, Collect}
Favourably grant, we beseech thee, O Lord, thy servants both pardon and peace;
that, being cleansed from the guilt of all their offences, they may serve thee
with secure minds. Thro'.
\end{englishwitness}

\noindent ``Favourably'' stands for \latin{placatus}, which says that God
grants as one appeased. ``Thy servants'' stands for \latin{fidelibus tuis}, thy
faithful. The English does not render \latin{pariter}, alike or together, and
``the guilt of'' is the translator's addition. The printed conclusion is short and is said in the full
form of the general rubrics (RG 115~a). In a low or conventual Mass on
11~October 2026 the Collect of the Maternity follows under its own conclusion
(below).

\subsection{Epistle}
\label{proper-epistle}

\properlocus{\latin{Lectio Epistolae beati Pauli Apostoli ad Ephesios}. Eph 5:15--21. Missal pp.~412--413, no.~1691.}

\begin{latinproper}
Fratres: Vidéte quómodo caute ambulétis: non quasi insipiéntes, sed ut
sapiéntes, rediméntes tempus, quóniam dies mali sunt. Proptérea nolíte fíeri
imprudéntes, sed intellegéntes quæ sit volúntas Dei. Et nolíte inebriári vino,
in quo est luxúria: sed implémini Spíritu Sancto, loquéntes vobismetípsis in
psalmis, et hymnis, et cánticis spirituálibus, cantántes, et psalléntes in
córdibus vestris Dómino: grátias agéntes semper pro ómnibus, in nómine Dómini
nostri Iesu Christi, Deo et Patri. Subiécti ínvicem in timóre Christi.
\end{latinproper}

\begin{englishwitness}{Douay--Rheims (Challoner), Eph 5:15--21}
See therefore, brethren, how you walk circumspectly: not as unwise, But as wise:
redeeming the time, because the days are evil. Wherefore, become not unwise:
but understanding what is the will of God. And be not drunk with wine, wherein
is luxury: but be ye filled with the Holy Spirit, Speaking to yourselves in
psalms and hymns and spiritual canticles, singing and making melody in your
hearts to the Lord: Giving thanks always for all things, in the name of our Lord
Jesus Christ, to God and the Father: Being subject one to another, in the fear
of Christ.
\end{englishwitness}

\noindent The Missal opens the lesson with the address \latin{Fratres} and
leaves out the Vulgate's connective and vocative, \latin{Videte itaque,
fratres}, which the English renders ``See therefore, brethren''. It spells
\latin{intellegentes} where the Clementine has \latin{intelligentes}. Otherwise
the Latin agrees with the Vulgate word for word.

\subsection{Gradual}
\label{proper-gradual}

\properlocus{Ps 144:15--16 (Heb.\ 145), respond and verse. Missal p.~413, no.~1692.}

\begin{latinproper}
Oculi ómnium in te sperant, Dómine: et tu das illis escam in témpore
opportúno. ℣.~Aperis tu manum tuam: et imples omne ánimal benedictióne.
\end{latinproper}

\begin{englishwitness}{Douay--Rheims (Challoner), Ps 144:15--16}
The eyes of all hope in thee, O Lord: and thou givest them meat in due season.
Thou openest thy hand, and fillest with blessing every living creature.
\end{englishwitness}

\noindent The Missal reads \latin{et tu das illis escam}, ``and thou givest them
food'', where the Clementine reads \latin{et tu das escam illorum}, ``their
food''; the English renders either. The verse agrees with the Vulgate.

\subsection{Alleluia}
\label{proper-alleluia}

\properlocus{Ps 107:2 (Heb.\ 108:2; English versions 108:1). Missal p.~413, no.~1693.}

\begin{latinproper}
Allelúia, allelúia. ℣.~\textnormal{Ps.\ 107, 2} Parátum cor meum, Deus,
parátum cor meum: cantábo, et psallam tibi, glória mea. Allelúia.
\end{latinproper}

\begin{englishwitness}{Douay--Rheims (Challoner), Ps 107:2}
My heart is ready, O God, my heart is ready: I will sing, and will give praise,
with my glory.
\end{englishwitness}

\noindent The English carries the first half of the verse and not the Missal's
close. The Clementine reads \latin{cantabo, et psallam in gloria mea}, which
the Douay renders; the Missal reads \latin{cantabo, et psallam tibi, gloria
mea}. It addresses the praise to God, \latin{tibi}, and sets \latin{gloria mea}
without its preposition, so that God himself is named the singer's glory. The
words \latin{psallam tibi} and \latin{gloria mea} both stand in the psalm, at
verses 4 and 3, but not together in a single verse. There is no Tract and no
Sequence.

\subsection{Gospel}
\label{proper-gospel}

\properlocus{\latin{Sequentia sancti Evangelii secundum Ioannem}. Jn 4:46--53. Missal p.~413, no.~1694.}

\begin{latinproper}
In illo témpore: Erat quidam régulus, cuius fílius infirmabátur Caphárnaum.
Hic cum audísset, quia Iesus adveníret a Iudǽa in Galilǽam, ábiit ad eum, et
rogábat eum ut descénderet, et sanáret fílium eius: incipiébat enim mori.
Dixit ergo Iesus ad eum: Nisi signa et prodígia vidéritis, non créditis.
Dicit ad eum régulus: Dómine, descénde priúsquam moriátur fílius meus. Dicit
ei Iesus: Vade, fílius tuus vivit. Crédidit homo sermóni, quem dixit ei Iesus,
et ibat. Iam autem eo descendénte, servi occurrérunt ei, et nuntiavérunt
dicéntes, quia fílius eius víveret. Interrogábat ergo horam ab eis, in qua
mélius habúerit. Et dixérunt ei: Quia heri hora séptima relíquit eum febris.
Cognóvit ergo pater quia illa hora erat, in qua dixit ei Iesus: Fílius tuus
vivit: et crédidit ipse, et domus eius tota.
\end{latinproper}

\begin{englishwitness}{Douay--Rheims (Challoner), Jn 4:46--53, from the middle of verse 46}
And there was a certain ruler, whose son was sick at Capharnaum. He having heard
that Jesus was come from Judea into Galilee, went to him and prayed him to come
down and heal his son: for he was at the point of death. Jesus therefore said to
him: Unless you see signs and wonders, you believe not. The ruler saith to him:
Lord, come down before that my son die. Jesus saith to him: Go thy way. Thy son
liveth. The man believed the word which Jesus said to him and went his way. And
as he was going down, his servants met him: and they brought word, saying, that
his son lived. He asked therefore of them the hour wherein he grew better. And
they said to him: Yesterday at the seventh hour, the fever left him. The father
therefore knew that it was at the same hour that Jesus said to him: Thy son
liveth. And himself believed, and his whole house.
\end{englishwitness}

\noindent The Missal's opening, \latin{In illo tempore: Erat quidam regulus},
stands in place of the first half of verse 46, ``He came again therefore into
Cana of Galilee, where he made the water wine'', so that the lesson begins
inside the verse it cites and does not name Cana; the English above begins
where the Missal's words rejoin the Vulgate. Every other word agrees with the
Vulgate. The printed cue \latin{Credo} follows the last line.

\subsection{Offertory}
\label{proper-offertory}

\properlocus{Ps 136:1 (Heb.\ 137). Missal p.~413, no.~1695.}

\begin{latinproper}
Super flúmina Babylónis illic sédimus, et flévimus: dum recordarémur tui,
Sion.
\end{latinproper}

\begin{englishwitness}{Douay--Rheims (Challoner), Ps 136:1}
Upon the rivers of Babylon, there we sat and wept: when we remembered Sion:
\end{englishwitness}

\noindent The Missal leaves out the psalm's Vulgate title, \latin{Psalmus David,
Ieremiae}, reads \latin{dum} for the Clementine's \latin{cum}, and adds
\latin{tui}: the captives remember Sion by speaking to her, ``thee, Sion'',
where the Vulgate and the English name her in the third person. The same words,
\latin{dum recordaremur tui, Sion}, stand in the psalm text that Cassiodorus
expounds.

\subsection{Secret}
\label{proper-secret}

\properlocus{\latin{Secreta}. Missal p.~413, no.~1696.}

\begin{latinproper}
Cæléstem nobis prǽbeant hæc mystéria, quǽsumus, Dómine, medicínam: et vítia
nostri cordis expúrgent. Per Dóminum.
\end{latinproper}

\begin{englishwitness}{Cummiskey, The Roman Missal (1861), XX Sunday after Pentecost, Secret}
May these mysteries, O Lord, we beseech thee, procure us a heavenly remedy, and
cleanse away the vices of our hearts. Thro'.
\end{englishwitness}

\noindent ``Remedy'' stands for \latin{medicinam}, medicine; the mysteries are
the subject of both verbs, so that what is asked is that the gifts themselves
heal and purge. The Missal follows the Secret with the direction
\latin{Praefatio de Ss.ma Trinitate}.

\subsection{Communion}
\label{proper-communion}

\properlocus{Ps 118:49--50 (Heb.\ 119). Missal p.~413, no.~1697.}

\begin{latinproper}
Meménto verbi tui servo tuo, Dómine, in quo mihi spem dedísti: hæc me
consoláta est in humilitáte mea.
\end{latinproper}

\begin{englishwitness}{Douay--Rheims (Challoner), Ps 118:49--50, to the Missal's close}
Be thou mindful of thy word to thy servant, in which thou hast given me hope.
This hath comforted me in my humiliation:
\end{englishwitness}

\noindent The Missal reads \latin{Memento} for the Vulgate's \latin{Memor esto},
adds the vocative \latin{Domine} after \latin{servo tuo}, and stops before the
end of verse 50, which in the Vulgate gives the reason for the comfort,
\latin{quia eloquium tuum vivificavit me}, ``because thy word hath enlivened
me''. The Vulgate counts the letter \latin{Zain} inside verse 49; the Missal
leaves it out.

\subsection{Postcommunion}
\label{proper-postcommunion}

\properlocus{\latin{Postcommunio}. Missal p.~413, no.~1698.}

\begin{latinproper}
Ut sacris, Dómine, reddámur digni munéribus: fac nos, quǽsumus, tuis semper
obœdíre mandátis. Per Dóminum nostrum.
\end{latinproper}

\begin{englishwitness}{Cummiskey, The Roman Missal (1861), XX Sunday after Pentecost, Postcommunion}
That we may be worthy of thy sacred gifts, O Lord: grant, we beseech thee, we
may always obey thy commandments. Thro'.
\end{englishwitness}

\noindent \latin{Reddamur} is passive, ``that we may be made worthy''; the
English's ``that we may be worthy'' leaves the making unnamed. The noun and verb
of the petition, \latin{mandatis} and \latin{oboedire}, are those of the
Introit's confession, \latin{mandatis tuis non oboedivimus}. The Missal prints
the conclusion \latin{Per Dominum nostrum}, longer than the cue on the Collect
and Secret.

\subsection{The commemoration in a low Mass on 11 October 2026}

These three orations belong to the formulary of the Maternity of the Blessed
Virgin Mary (\latin{Die 11 octobris}, Missal pp.~685--686, nos.~3847, 3855 and
3857), not to the Sunday. In a low Mass, and in a conventual Mass, each follows
the Sunday's oration of the same kind under its own conclusion; a sung Mass that
is not conventual omits them. Their wording, too, stands in the 1862 Pustet
Missal, among its Masses for certain places (appendix pp.~[171]--[172]). The
English is that of F.~X. Lasance's \work{The New Roman Missal}, in the
Benziger revision of 1945.

\begin{latinproper}
Deus, qui de beátæ Maríæ Vírginis útero Verbum tuum, Angelo nuntiánte, carnem
suscípere voluísti: præsta supplícibus tuis; ut, qui vere eam Genetrícem Dei
crédimus, eius apud te intercessiónibus adiuvémur. Per eúndem Dóminum nostrum
Iesum Christum Fílium tuum: Qui tecum vivit et regnat in unitáte.
\end{latinproper}

\begin{englishwitness}{Lasance, The New Roman Missal (1945), Maternity of the Blessed Virgin Mary, Collect}
O God, Who wast pleased that at the angel's message Thy Word should take flesh
in the womb of the blessed Virgin Mary, grant to Thy suppliants that, believing
her to be truly the mother of God, we may be assisted by her intercessions with
Thee. Through the same.
\end{englishwitness}

\begin{latinproper}
Tua, Dómine, propitiatióne, et beátæ Maríæ semper Vírginis, Unigéniti tui
Matris, intercessióne, ad perpétuam atque præséntem hæc oblátio nobis
profíciat prosperitátem et pacem. Per eúndem Dóminum.
\end{latinproper}

\begin{englishwitness}{Lasance, The New Roman Missal (1945), Maternity of the Blessed Virgin Mary, Secret}
Through Thy mercy, O Lord, and the intercession of blessed Mary, ever a virgin,
the Mother of Thine only-begotten Son, may our oblation profit us for eternal
and for present prosperity and peace. Through the same.
\end{englishwitness}

\begin{latinproper}
Hæc nos commúnio, Dómine, purget a crímine: et, intercedénte beáta Vírgine Dei
Genetríce María, cæléstis remédii fáciat esse consórtes. Per eúndem Dóminum.
\end{latinproper}

\begin{englishwitness}{Lasance, The New Roman Missal (1945), Maternity of the Blessed Virgin Mary, Postcommunion}
May this communion, O Lord, purge away our guilt and, by the intercession of
blessed Mary the Mother of God, make us companions of Him, Who is our heavenly
healing. Through the same.
\end{englishwitness}

\noindent The Postcommunion's Latin asks that the communion make us
\latin{consortes}, sharers, \latin{caelestis remedii}, of the heavenly remedy;
the 1945 translators read that remedy as Christ himself. In a low Mass the two
Secrets therefore stand side by side, the Sunday's asking the mysteries to be
heavenly medicine and the Maternity's asking prosperity and peace, and the two
Postcommunions ask, the one obedience to the commandments, the other cleansing
from guilt and a share in the heavenly remedy.
